%
%============= CHAPTER BREAK =================%
%
\chapter{Philostratus II (213--217 CE)}
\vspace*{3mm}%
\justifying%
%
\noindent{}Philostratus II the Sophist % 
is best known in early Christian studies for his %
\textit{Life of Apollonius of Tyana} / \textit{Vita Apollonii}, %
which is often seen as a satirical rival to early Christian Gospels. %
But his \textit{Heroikos}, written between 213--217 CE,\endnote{%
Jones 2001:143.} %
has also received some attention by scholars of early Christianity.\endnote{%
See esp. Berenson Maclean \& Aitken 2001.} 
He lived and studied in Athens, taught in elite circles in Rome, %
and eventually moved to Tyre (modern Lebanon), where he died. \\[1mm] % %

\noindent{}\textbf{ἀμαράντινος} / \textit{amarantinos} / “amaranthine” in TLG results has two %
candidates for the earliest attestation among non-Christian sources. %
The first of these is in Ps-Galen, \textit{De succedaneis}, or \textit{On Substitutions}, %
a text the TLG notes is after the 2nd century CE: %
“instead of sesamoid, \textbf{amaranthine} pulp” / ἀντὶ σησαμοειδοῦς, πίεσμα \textbf{ἀμαραντινόν}.\endnote{%
Kühn 1830 19:742.}

The most chronologically secure instance, despite the difficulties in dating the work precisely,\endnote{%
Berenson Maclean \& Aitken 2001:xlii--xlv.} %
is in Philostratus II, \textit{Heroikos} 53.9: %

\vspace*{3mm} 
 \begin{quote}κατ’ ἀρχὰς μὲν δὴ τοιάδε ἐγίνετο· %
 ναῦς ἐκ Θετταλίας μέλανα ἱστία ἠρμένη ἐς Τροίαν ἔπλει, %
 θεωροὺς μὲν δὶς ἑπτὰ ἀπάγουσα, ταύρους δὲ λευκόν τε καὶ μέλανα, %
 χειροήθεις ἄμφω, καὶ ὕλην ἐκ Πηλίου, %
 ὡς μηδὲν τῆς πόλεως δέοιντο· %
 καὶ πῦρ ἐκ Θετταλίας ἦγον καὶ σπονδὰς καὶ ὕδωρ τοῦ Σπερχειοῦ ἀρυσάμενοι· %
 ὅθεν καὶ στεφάνους \textbf{ἀμαραντίνους} ἐς τὰ κήδη πρῶτοι Θετταλοὶ ἐνόμισαν, %
 ἵνα, κἂν ἄνεμοι τὴν ναῦν ἀπολάβωσι, μὴ σαπροὺς ἐπιφέρωσι μηδὲ ἐξώρους. \\[1mm] %
 To begin with, these things indeed happened: %
 a ship from Thessaly, raising black sails, was sailing to Troy, %
 carrying off twice seven spectators, and bulls both white and black, %
 both hand-tamed, and wood from Pelion, %
 so they would need nothing from the city. %
 Fire from Thessaly they also brought, and drink-offerings and %
 water drawn from the Spercheios. %
 Wherefore the Thessalians first ordained \textbf{amaranthine} crowns for funeral rites, %
 so that, if winds should take away the ship, they would not carry on rotted or too old.\endnote{%
 Berenson Maclean \& Aitken translate ἀμαραντίνους as “unfading”. %
 Our translation is more literal, somewhat akin to the custom in English %
 to speak of an “evergreen” wreath or tree. For this naturalistic usage, %
 “unfading” is not inaccurate, and yet is too generic.} %
\end{quote} % 
\vspace*{3mm}

In its context, this describes offerings by Thessalians made to the hero Achilles, %
in conformity with the oracle of Dodona, %
to travel to Troy to make sacrifices to Achilles and to the dead (\textit{Heroik}. 53.8).
The term appears only once in the NT, specifically in %
a passage from I Peter with guidance for community leaders (1Pet 5:2--4): %

\vspace*{3mm} 
 \begin{quote}ποιμάνατε τὸ ἐν ὑμῖν ποίμνιον τοῦ θεοῦ, %
 ἐπισκοποῦντες μὴ ἀναγκαστῶς ἀλλὰ ἑκουσίως κατὰ θεόν, %
 μηδὲ αἰσχροκερδῶς ἀλλὰ προθύμως, 
 ³ μηδ’ ὡς κατακυριεύοντες τῶν κλήρων ἀλλὰ τύποι γινόμενοι τοῦ ποιμνίου· %
 ⁴ καὶ φανερωθέντος τοῦ ἀρχιποίμενος κομιεῖσθε τὸν \textbf{ἀμαράντινον} τῆς δόξης στέφανον. \\[1mm] %
 Shepherd the flock of the God among you, %
 over-watching not out of necessity but instead willingly according to God, %
 neither sordidly greedy, but instead devotedly, %
 ³ not as those who over-lord the allotments, but instead becoming models for the flock. %
 ⁴ And when the chief-shepherd is manifested, you will inherit the \textbf{amaranthine} crown of glory. \end{quote} % 
 \vspace*{3mm}

Even beyond the rare word, the similarities are striking. %
Both use the new/rare term in connection with a crown. %
But are addressed to leaders of the people in regard to their religious responsibilities. %
Once these connections are clarified, %
other elements of I Peter here may hint at the role of Achilles as %
one who (unlike Agamemnon), did not lord it over his charges, %
but instead became a model of courage to them. %
The theme of kings as shepherds is commonplace in Homeric lore; %
one might say that Philostratus pictures Achilles as a model “chief-shepherd” who dies for the flock, %
one who receives an undying crown of eternal fame and veneration. %
It thus also locates the emerging cult of Jesus Christ squarely within %
traditional literary practices of Greco-Roman hero cult. %
Note also that ἀρχιποίμην / \textit{archipoimēn} / “chief-shepherd” was already identified as an Aurelian-era neologism %
in the chapter on Herodian the Grammarian. %
Thus in this verse in I Peter we find two neologisms in a single verse, %
mutually strengthening the case for late redaction. %
Though I Peter is securely attested in large part by Irenaeus ca 180 CE, %
these findings demonstrate it was still being edited well into the early 3rd century CE. %
